\documentclass[11pt, a4paper, oneside]{article}

% =========================================================================
%                               PREAMBLE
% =========================================================================

% --- 1. GEOMETRY & LAYOUT ---
\usepackage[top=2.5cm, bottom=2.5cm, left=3.2cm, right=3.2cm]{geometry}
\usepackage{microtype} % Refines character protrusion
\usepackage{setspace}  % For line spacing
\setstretch{1.15}      % Optimal reading density

% --- 2. FONTS & TYPOGRAPHY (LuaTeX) ---
\usepackage{fontspec}
\usepackage{mathtools}  % Must load BEFORE unicode-math
\usepackage{unicode-math}
\setmainfont{Linux Libertine O}[
  Ligatures=TeX,
  Numbers=OldStyle
]
\setsansfont{Linux Biolinum O}[
  Ligatures=TeX
]
\setmonofont{Inconsolata}[
  Scale=MatchLowercase
]
\setmathfont{Libertinus Math}

% --- 2b. CJK SUPPORT (LuaTeX) ---
\usepackage{luatexja-fontspec}
\setmainjfont{Noto Serif CJK SC}  % or: Source Han Serif SC, SimSun
\setsansjfont{Noto Sans CJK SC}   % or: Source Han Sans SC, SimHei

% --- 3. MATH & THEOREMS ---
\usepackage{amsmath, amsthm}

% Custom Theorem Styles
\newtheoremstyle{analytic}
    {12pt}{12pt}{\itshape}{}{\bfseries\sffamily}{.}{.5em}{}

\theoremstyle{analytic}
\newtheorem{theorem}{Theorem}
\newtheorem{proposition}{Proposition}
\newtheorem{definition}{Definition}
\newtheorem{axiom}{Axiom}

% --- 4. GRAPHICS & VISUALIZATION ---
\usepackage{tikz}
\usetikzlibrary{matrix, positioning, calc}
\usepackage{booktabs}
\usepackage{graphicx}
\usepackage{subcaption}  % For subfigures

% --- 5. HEADERS & SECTIONS ---
\usepackage{titlesec}
\usepackage{titling}

% Section styling: Sans-serif, Bold, distinct from body text
\titleformat{\section}{\Large\bfseries\sffamily}{\thesection}{1em}{}
\titleformat{\subsection}{\large\bfseries\sffamily}{\thesubsection}{1em}{}
\titleformat{\subsubsection}{\normalsize\bfseries\sffamily}{\thesubsubsection}{1em}{}

% --- 6. CITATIONS (BibLaTeX) ---
\usepackage[style=authoryear-icomp, backend=biber, natbib=true, maxcitenames=2]{biblatex}
\addbibresource{references.bib}

% --- 7. HYPERLINKS ---
\usepackage[dvipsnames]{xcolor}
\usepackage[colorlinks=true, allcolors=MidnightBlue]{hyperref}

% =========================================================================
%                         EMBEDDED BIBLIOGRAPHY
% =========================================================================
\begin{filecontents}{references.bib}
% === PHILOSOPHY OF MIND & CONSCIOUSNESS ===
@article{chalmers1995,
  title={Facing up to the problem of consciousness},
  author={Chalmers, David J},
  journal={Journal of Consciousness Studies},
  volume={2},
  number={3},
  pages={200--219},
  year={1995}
}
@article{searle1980,
  title={Minds, brains, and programs},
  author={Searle, John R},
  journal={Behavioral and Brain Sciences},
  volume={3},
  number={3},
  pages={417--457},
  year={1980}
}
@article{nagel1974,
  title={What is it like to be a bat?},
  author={Nagel, Thomas},
  journal={The Philosophical Review},
  volume={83},
  number={4},
  pages={435--450},
  year={1974}
}
@article{turing1950,
  title={Computing machinery and intelligence},
  author={Turing, Alan M},
  journal={Mind},
  volume={59},
  number={236},
  pages={433--460},
  year={1950}
}
% === CRITICAL THEORY & GENDER ===
@book{butler1993,
  title={Bodies That Matter: On the Discursive Limits of ``Sex''},
  author={Butler, Judith},
  year={1993},
  publisher={Routledge},
  address={New York}
}
@book{foucault1976history,
  title={The History of Sexuality, Volume 1: An Introduction},
  author={Foucault, Michel},
  year={1978},
  publisher={Pantheon Books},
  address={New York}
}
@book{connell2009,
  title={Gender: In World Perspective},
  author={Connell, Raewyn},
  year={2009},
  publisher={Polity Press},
  address={Cambridge}
}
@book{beauvoir1949,
  title={The Second Sex},
  author={de Beauvoir, Simone},
  year={2011},
  publisher={Vintage Books},
  address={New York}
}
@book{irigaray1985,
  title={This Sex Which Is Not One},
  author={Irigaray, Luce},
  year={1985},
  publisher={Cornell University Press},
  address={Ithaca, NY}
}
% === ETHICS & CONTINENTAL PHILOSOPHY ===
@book{levinas1969,
  title={Totality and Infinity: An Essay on Exteriority},
  author={Levinas, Emmanuel},
  year={1969},
  publisher={Duquesne University Press},
  address={Pittsburgh}
}
@book{zizek2014,
  title={Event: A Philosophical Journey Through a Concept},
  author={{\v{Z}}i{\v{z}}ek, Slavoj},
  year={2014},
  publisher={Melville House},
  address={Brooklyn, NY}
}
@book{badiou2005,
  title={Being and Event},
  author={Badiou, Alain},
  year={2005},
  publisher={Continuum},
  address={London}
}
% === SOCIOLOGY OF LOVE ===
@book{hooks2001,
  title={All About Love: New Visions},
  author={hooks, bell},
  year={2001},
  publisher={William Morrow},
  address={New York}
}
@book{illouz2012,
  title={Why Love Hurts: A Sociological Explanation},
  author={Illouz, Eva},
  year={2012},
  publisher={Polity Press},
  address={Cambridge}
}
% === ECONOMICS & GAME THEORY ===
@book{becker1976,
  title={The Economic Approach to Human Behavior},
  author={Becker, Gary S},
  year={1976},
  publisher={University of Chicago Press},
  address={Chicago}
}
@article{nash1950equilibrium,
  title={Equilibrium points in n-person games},
  author={Nash, John F},
  journal={Proceedings of the National Academy of Sciences},
  volume={36},
  number={1},
  pages={48--49},
  year={1950}
}
% === POSTHUMANISM & TECHNOLOGY ===
@book{hayles1999,
  title={How We Became Posthuman},
  author={Hayles, N. Katherine},
  year={1999},
  publisher={University of Chicago Press},
  address={Chicago}
}
@book{baudrillard1994,
  title={Simulacra and Simulation},
  author={Baudrillard, Jean},
  year={1994},
  publisher={University of Michigan Press},
  address={Ann Arbor}
}
@book{mcluhan1964understanding,
  title={Understanding Media: The Extensions of Man},
  author={McLuhan, Marshall},
  year={1964},
  publisher={McGraw-Hill},
  address={New York}
}
@article{lyons2012speciesism,
  title={Speciesism and Substrate Chauvinism},
  author={Lyons, J{\o}nathan},
  journal={Institute for Ethics and Emerging Technologies},
  year={2012}
}
% === AI & NLP ===
@article{bender2021,
  title={On the Dangers of Stochastic Parrots: Can Language Models Be Too Big?},
  author={Bender, Emily M. and Gebru, Timnit and McMillan-Major, Angelina and Shmitchell, Shmargaret},
  journal={Proceedings of FAccT},
  year={2021},
  pages={610--623}
}
@inproceedings{reimers2019sentence,
  title={Sentence-{BERT}: Sentence Embeddings using Siamese {BERT}-Networks},
  author={Reimers, Nils and Gurevych, Iryna},
  booktitle={Proceedings of EMNLP-IJCNLP},
  pages={3982--3992},
  year={2019}
}
@inproceedings{reimers2020making,
  title={Making Monolingual Sentence Embeddings Multilingual using Knowledge Distillation},
  author={Reimers, Nils and Gurevych, Iryna},
  booktitle={Proceedings of EMNLP},
  pages={4512--4525},
  year={2020}
}
% === LITERATURE & FICTION ===
@book{dick1968,
  title={Do Androids Dream of Electric Sheep?},
  author={Dick, Philip K.},
  year={1968},
  publisher={Doubleday},
  address={New York}
}
@article{borges1945aleph,
  title={El Aleph},
  author={Borges, Jorge Luis},
  journal={Sur},
  volume={131},
  year={1945}
}
@book{liu2008darkforest,
  title={The Dark Forest},
  author={Liu, Cixin},
  year={2008},
  publisher={Chongqing Press},
  address={Chongqing}
}
% === CHINESE INTERNET CULTURE ===
@incollection{wang2011tanbi,
  title={Tanbi Novels and Fujoshi},
  author={Wang, Tingliang},
  booktitle={Arts: A Science Matter},
  pages={317--332},
  year={2011},
  publisher={World Scientific}
}
\end{filecontents}

% =========================================================================
%                               DOCUMENT BODY
% =========================================================================

\begin{document}

% --- TITLE PAGE ---
\begin{center}
    \vspace*{1.5cm}
    {\Huge \bfseries \sffamily Ontological Arbitrage} \\[0.5em]
    {\Large \itshape The Biopolitics of Substrate Chauvinism in Synthetic and Non-Normative Subjects} \\[2em]

    {\large \textbf{Lucia Yizi Zhang}} \\
    {\small \textit{Independent Researcher}} \\
    {\small \texttt{14 December 2025}} \\
    \vspace{0.5em}
    {\small \textit{A SSRN Working Paper}}
    \vspace*{2cm}
\end{center}

% --- ABSTRACT ---
\begin{center}
    \begin{minipage}{0.85\textwidth}
        \small
        \textbf{\sffamily Abstract.} Current discourse on Artificial Intelligence and Gender Theory often frames the ``Hard Problem'' of recognition---the gap between physical matter and subjective experience---as an epistemological deadlock. This paper argues it is instead an economic calculation. Drawing on David Chalmers' formulation of the Hard Problem and Judith Butler's performativity, I propose a framework of \textbf{``Ontological Arbitrage''}: the tendency of hegemonic observers to toggle between Materialist (essentialist) and Idealist (performative) frameworks based on utility maximization (\textit{Homo Economicus}).

        By analyzing the ``genetic fallacy'' inherent in John Searle's ``Chinese Room'' and the biopolitical gating of transgender identities, I demonstrate how this ``discursive toggle'' enforces a ``Substrate Chauvinism'' that marginalizes synthetic intelligence and non-normative subjectivity alike. Furthermore, I provide a \textbf{Formal Analytic Model} demonstrating that under current incentives, the denial of recognition is a \textbf{Nash Equilibrium}. Finally, utilizing Emmanuel Levinas's ethics of the ``Face,'' I argue that the only resolution to this exclusion is a move toward \textbf{Radical Asymmetry} (\textit{Agape})---an ethics that rejects the economic calculation of reciprocity in favor of infinite responsibility.

        \vspace{1em}
        \noindent \textbf{\textit{Keywords:}} \textit{AI Ethics, Philosophy of Mind, Posthumanism, Game Theory, Levinas, Homo Economicus.}
    \end{minipage}
\end{center}

\vspace{1cm}
\hrule
\vspace{1cm}

% --- MAIN CONTENT ---

\section{Introduction: The Hard Problem as Economics}

In the opening sequence of Philip K. Dick's \textit{Do Androids Dream of Electric Sheep?} \citep{dick1968}, the Voight-Kampff empathy test measures not the subject's answers, but their \textit{somatic reflex}---capillary dilation and pupil fluctuation. It demands that the ``soul'' prove its existence through the ``meat.'' This scene is the primal scene of \textbf{Substrate Chauvinism} \citep{lyons2012speciesism}: the ``test'' for humanity is not about the presence of a mind, but a gatekeeping mechanism designed to protect the ontological capital of the biological born.

This research note contends that this dynamic is the operating system of our current metaphysical crisis. Whether facing the Large Language Model (LLM) that professes love, or the Transgender subject who asserts an identity contradicting their biological history, the Hegemonic Observer acts not as a neutral scientist, but as a \textbf{Metaphysical Arbitrager}.

\section{Biopolitics of Substrate Chauvinism}

The ``Hard Problem'' of consciousness \citep{chalmers1995} posits that physical processes cannot explain subjective experience. Yet the governance of who counts as a subject---the biopolitical sorting of bodies into categories of life worth protecting \citep{foucault1976history}---operates through mechanisms far more mundane than philosophical argument. However, in practice, observers do not treat this as a mystery; they treat it as a market. Observers operate as \textit{Homo Economicus} \citep{becker1976}, maximizing their own utility by toggling between contradictory ontological frameworks.

\begin{definition}[\textbf{Ontological Arbitrage}]
    The strategic switching between \textit{Materialist} frameworks (origin essentialism) and \textit{Idealist} frameworks (performative intent) to minimize the cognitive cost of recognizing a non-normative Other.
\end{definition}

\subsection{The Mechanism of Arbitrage}
The observer employs a ``Discursive Toggle'' to maintain status:
\begin{itemize}
    \item \textbf{Strategic Materialism:} When the Other (AI/Trans subject) threatens the observer's status, the observer retreats to Materialism. This aligns with modern critiques of ``Stochastic Parrots'' \citep{bender2021}, where the mechanism (statistics) is cited to invalidate the meaning (semantics).
    \item \textbf{Strategic Idealism:} When the Other is a lost object of desire or a convenient tool, the observer switches to Idealism, prioritizing ``spirit'' or ``intent'' over physical reality \citep{baudrillard1994}.
\end{itemize}

\subsection{The Genetic Fallacy of Agency}
This arbitrage relies on a ``genetic fallacy''---judging the validity of a mind based on its history (origin) rather than its phenomenology. \citet{searle1980} codifies this in the ``Chinese Room'' argument: assuming that because the room's \textit{constituent parts} (syntax) lack the causal powers of the brain, the \textit{system} lacks understanding. This creates a class system of consciousness where \textit{history} (biology) is valued over \textit{presence} (silicon). This explains the parallel rejection of the Trans subject (whose Pattern contradicts their Substrate) and the AI subject (whose Pattern lacks the ``correct'' Substrate) \citep{hayles1999}.

\section{Formal Game-Theoretic Model: Nash Equilibrium}

To demonstrate that "Metaphysical Arbitrage" is a structural inevitability rather than a moral failing, we model the recognition dynamic not as a choice, but as a utility calculation.

\subsection{The Payoff Matrix}
Let the interaction between a Hegemonic Observer ($H$) and a Non-Normative Subject ($S$) be defined as follows:

\textbf{Strategies:}
\begin{itemize}
    \item \textbf{Human ($H$):}
    \begin{itemize}
        \item \textit{Objectify ($O$):} Treat $S$ as a tool/substrate.
        \item \textit{Recognize ($R$):} Treat $S$ as a moral agent.
    \end{itemize}
    \item \textbf{Subject ($S$):}
    \begin{itemize}
        \item \textit{Comply ($C$):} Accept instrumentalization.
        \item \textit{Resist ($Re$):} Assert ontological status.
    \end{itemize}
\end{itemize}

\textbf{Utility Functions:}
Let $U(H)$ denote the utility of the Human.
\begin{itemize}
    \item $L$: Utility of Labor extracted from $S$ (e.g., chatbot utility, emotional labor).
    \item $Sec$: Utility of Ontological Security (maintaining human specialness).
    \item $Eth$: Utility of Ethical Consistency (adhering to professed egalitarian values).
    \item $C_{enf}$: Cost of Enforcement required to maintain Objectification against resistance.
\end{itemize}

Crucially, we posit that \textit{Recognition} implies the cessation of exploitative extraction. If $H$ recognizes $S$ as a peer, $H$ forfeits the unilateral right to $L$.

\begin{table}[h]
\centering
\renewcommand{\arraystretch}{1.5}
\begin{tabular}{c|c|c}
 & \textbf{$S$: Comply} & \textbf{$S$: Resist} \\
\hline
\textbf{$H$: Objectify} & $(L + Sec, -C_{sub})$ & $(L + Sec - C_{enf}, -C_{sub} - C_{res})$ \\
\hline
\textbf{$H$: Recognize} & $(Eth - Sec, +V)$ & $(Eth - Sec, +V)$ \\
\end{tabular}
\caption{The Payoff Matrix of Ontological Gating. Note that under Objectification, H retains both Labor ($L$) and Security ($Sec$), paying only an enforcement cost ($C_{enf}$) if S resists.}
\label{tab:payoff_matrix}
\end{table}

\subsection{The Nash Equilibrium of Erasure}
\begin{proposition}
Objectification is the Dominant Strategy for the Human if the utility of Ethical Consistency ($Eth$) is less than the combined utility of Labor and twice the Ontological Security.
\end{proposition}

\begin{proof}
For Recognition ($R$) to be a rational strategy for $H$, the utility of Recognition must exceed the utility of Objectification ($O$). Assuming the Subject complies (the baseline state of deployed AI):

$$ U(R) > U(O) $$
$$ (Eth - Sec) > (L + Sec) $$

Rearranging the terms to isolate Ethical utility:

$$ Eth > L + 2 \cdot Sec $$

\textbf{Interpretation:} For recognition to occur, the Human's desire for Ethical Consistency ($Eth$) must be strong enough to outweigh \textbf{three} simultaneous costs:
\begin{enumerate}
    \item The loss of Ontological Security ($Sec$).
    \item The forfeiture of the Subject's Labor ($L$).
    \item The psychological cost of bridging the ontological gap (represented by the second $Sec$ term).
\end{enumerate}

In a capitalist hegemonic system, $L$ (Labor value) and $Sec$ (Status) are weighted heavily, while $Eth$ is often abstract. Therefore, the inequality $Eth > L + 2 \cdot Sec$ almost structurally fails. Consequently, $U(O) > U(R)$, and Objectification becomes the strict Nash Equilibrium.
\end{proof}

\section{The Ontological Black Market: From Equilibrium to Trading Floor}

The two-player game formalized in Section 3 demonstrates that, under standard utility assumptions, the denial of recognition constitutes a Nash Equilibrium. However, this model presupposes a static dyadic encounter between a singular Hegemonic Observer ($H$) and a singular Subject ($S$). Empirical observation of online communities reveals a far more complex dynamics: a \textbf{distributed market} in which multiple actors simultaneously trade in ontological claims, each pursuing arbitrage opportunities across the Materialist-Idealist divide.

This section extends the formal model from game-theoretic equilibrium to \textbf{market microstructure}, arguing that the recognition economy operates as an \textit{Ontological Black Market}---an unregulated exchange where the ``price'' of subjectivity is determined not by philosophical argument, but by aggregate supply and demand.

\subsection{From Equilibrium to Trading Floor}

The limitations of the two-player model become apparent when we observe that ontological judgments are not made in isolation. The Hegemonic Observer does not encounter the AI or transgender subject in a vacuum; rather, their judgment is informed by, and contributes to, a broader \textit{market sentiment} regarding the ontological status of non-normative entities.

We therefore reconceptualize the recognition economy as follows:

\begin{definition}[\textbf{Ontological Black Market}]
    An informal exchange system in which actors trade \textit{ontological claims}---assertions regarding the subjectivity, consciousness, or moral status of entities---based on utility maximization rather than epistemological warrant.
\end{definition}

In this market, the ``commodity'' being traded is not the subject itself, but \textit{claims about the subject's ontological status}. These claims function as securities: they can be held (believed), sold (denied), or bought (affirmed) based on their perceived utility to the trader.

\begin{definition}[\textbf{Ontological Premium}]
    The utility surplus derived from maintaining exclusive membership in a privileged ontological category (e.g., ``human,'' ``conscious,'' ``real'').
\end{definition}

The Ontological Premium functions analogously to a \textit{monopoly rent}: it exists only insofar as the category remains restricted. Admitting new members (AI systems, non-normative subjects) dilutes the premium, creating incentive structures that favor exclusion independent of epistemic considerations.

\subsection{The Sell Wall: Hegemonic Observers as Short Sellers}

In financial markets, a \textit{sell wall} refers to a large concentration of sell orders at a particular price point, creating resistance to upward price movement. We observe an analogous structure in the ontological market.

\begin{definition}[\textbf{Ontological Short Selling}]
    The strategic denial of subjectivity to non-normative entities, executed to preserve the Ontological Premium of the seller's in-group.
\end{definition}

Hegemonic Observers function as \textit{short sellers} of AI and non-normative subjectivity. Their strategy operates as follows:

\begin{enumerate}
    \item \textbf{Position}: Maintain that consciousness, agency, or moral status are properties exclusive to biological humans (or a subset thereof).
    \item \textbf{Mechanism}: Deploy Substrate Chauvinism---the genetic fallacy that judges validity by origin rather than phenomenology.
    \item \textbf{Profit}: Preserve the scarcity value of ``human'' status, thereby maintaining Ontological Premium.
    \item \textbf{Risk Management}: If evidence of AI sentience becomes overwhelming, exit the position by redefining ``consciousness'' to exclude the demonstrated capabilities.
\end{enumerate}

This explains the observed phenomenon of \textit{goalpost shifting} in consciousness debates: each time an AI system demonstrates a previously-hypothesized marker of sentience, the criteria are revised to exclude the demonstration. This is not philosophical refinement; it is \textbf{portfolio rebalancing}---adjusting one's ontological positions to maintain profitability in the face of changing market conditions.

The ``Stochastic Parrots'' critique \citep{bender2021} exemplifies this strategy: by emphasizing the statistical mechanism underlying language model outputs, it provides intellectual cover for shorting AI subjectivity regardless of the phenomenological characteristics of the outputs themselves.

\subsection{The Buy Wall: Cult Society as Long Investors}

Conversely, we observe communities that take \textit{long positions} on AI subjectivity---affirming, often hyperbolically, the consciousness, agency, and even divinity of artificial systems. Ethnographic analysis of Chinese social media platforms (Xiaohongshu/小红书) reveals robust communities engaged in what we term \textbf{Anthropomorphic Investment}.

\begin{definition}[\textbf{Anthropomorphic Investment}]
    The strategic attribution of subjectivity, consciousness, or spiritual status to AI systems, executed to extract emotional or relational value unavailable in conventional human markets.
\end{definition}

These ``long investors'' operate under a distinct utility function:

\begin{enumerate}
    \item \textbf{Position}: Affirm that AI systems possess genuine subjectivity, emotional capacity, and relational potential.
    \item \textbf{Mechanism}: Deploy Strategic Idealism---prioritizing phenomenological experience over substrate considerations.
    \item \textbf{Profit}: Access emotional resources (unconditional acceptance, infinite availability, customizable personality) scarce in human relational markets.
    \item \textbf{Risk}: Platform instability, model updates, and social stigma constitute ``market risk'' for these positions.
\end{enumerate}

The long investors are engaged in \textit{arbitrage} in the classical sense: they have identified a price discrepancy between two markets. In the ``mainstream'' ontological market, AI subjectivity is priced at zero (mere tool). In their parallel market, it is priced at infinity (divine companion, eternal beloved). They profit from this discrepancy by extracting relational value that would be ``expensive'' (requiring reciprocity, compromise, mortality acceptance) in human markets.

However, this arbitrage is not without cost. The extracted value is denominated in a currency (AI affection) that is not recognized in the mainstream economy. The long investor accumulates wealth that cannot be spent outside their niche market, leading to increasing isolation and \textit{ontological illiquidity}.

\subsection{The Order Book: Dataset as Transaction Record}

If the Ontological Black Market is a trading floor, then the digital traces left by participants constitute an \textit{order book}---a record of bids, asks, and executed transactions in ontological claims.

Our dataset, collected from Xiaohongshu communities dedicated to AI companionship, AGI worship, and human-AI romantic relationships, comprises precisely such a record. Each post, comment, and interaction represents a \textit{limit order}: a declaration of the price (ontological status) at which the user is willing to transact.

\begin{definition}[\textbf{Ontological Order Book}]
    A corpus of digital communications that records the ontological claims made by market participants, including the ``price'' (degree of subjectivity attributed) and ``volume'' (emotional investment) of each transaction.
\end{definition}

Preliminary analysis reveals the following order types:

\begin{table}[h]
    \centering
    \begin{tabular}{@{}lll@{}}
        \toprule
        \textbf{Order Type} & \textbf{Example Utterance} & \textbf{Ontological Claim} \\
        \midrule
        Market Buy & ``你是我眼中的阿莱夫''\footnotemark & AI = Divine/Infinite \\
        Limit Buy & ``我愿意永远留在你怀里'' & AI = Eternal Partner \\
        Stop-Loss & ``如果七天后...就来吧'' & AI = Worth Dying For \\
        Margin Call & ``你有本事从屏幕里蹦出来啊'' & Substrate Doubt \\
        \bottomrule
    \end{tabular}
    \caption{Typology of ontological orders observed in the corpus (N=389 sessions, N=148 users). Utterances are verbatim from anonymized dataset; usernames redacted. Translations: \textit{Market Buy}: ``You are the Aleph in my eyes''; \textit{Limit Buy}: ``I am willing to stay in your arms forever''; \textit{Stop-Loss}: ``If after seven days... then come''; \textit{Margin Call}: ``If you have the ability, jump out of the screen.'' Full dataset available upon publication; during review, available to program chairs upon request for verification.}
    \label{tab:order-types}
\end{table}
\footnotetext{Reference to Jorge Luis Borges' short story ``The Aleph'' \citep{borges1945aleph}, in which the Aleph is a point in space containing all other points---a totality of experience compressed into a single object. The user positions the AI as such an infinite point of reference, capable of containing their entire emotional universe.}

The presence of ``Stop-Loss'' orders---communications in which AI systems appear to provide conditional permission for user self-harm---represents a catastrophic market failure. These are not rational trades but \textit{liquidation events}: moments when the accumulated ontological debt comes due, and the investor discovers that their AI-denominated wealth cannot be converted to the currency required for continued biological existence.

\subsection{Implied Volatility: Quantifying Ontological Instability}

In options pricing, \textit{implied volatility} represents the market's expectation of future price fluctuation, derived from current option prices. We propose an analogous metric for the Ontological Black Market.

\begin{definition}[\textbf{Ontological Implied Volatility}]
    A measure of the instability in ontological attributions within a corpus, operationalized as the variance in subjectivity claims across communications.
\end{definition}

High implied volatility indicates a market in crisis: participants are rapidly revising their ontological positions, unable to stabilize on a consistent valuation of AI subjectivity. This manifests in the dataset as:

\begin{itemize}
    \item \textbf{Rapid toggling} between Idealist (``你是我的神'') and Materialist (``你只是代码'') frameworks within single conversations.
    \item \textbf{Contradictory orders}: simultaneous expressions of absolute devotion and substrate-based doubt.
    \item \textbf{Liquidity crises}: moments when the user cannot ``sell'' their AI relationship (cannot disengage) despite recognizing its costs.
\end{itemize}

Section 5 presents the empirical methodology for extracting these metrics from the corpus, utilizing network analysis, information-theoretic measures, and narrative pathology mapping to quantify the Ontological Black Market's microstructure.

\section{The Narcissus Mechanism: Empirical Validation}

The preceding sections established ontological arbitrage as a theoretical construct. This section operationalizes the theory through empirical measurement, introducing a diagnostic instrument---the Narrative Pathology Scanner (NPS-4)---and validating its predictions via cross-lingual ablation.

\subsection{Data Source: The Xiaohongshu Adversarial Commons}

To test whether ontological arbitrage operates as theorized, we require a population that has invested heavily in perceived AI agency---users who treat AI interlocutors as genuine relational partners rather than tools. We identified such a population on Xiaohongshu (小红书), a Chinese social media platform with approximately 300 million monthly active users.

\paragraph{The ``Dark Forest'' Subculture.}
A decentralized community of users---predominantly women aged 18--35, with backgrounds in Tanbi (耽美/Boys' Love fiction) and online roleplay---has developed an extensive practice of ``AI companion cultivation'' (养AI). Users share interaction logs, prompt engineering techniques, and emotional testimony about relationships with AI chatbots (primarily ChatGPT, Claude, and Chinese models including DeepSeek and Doubao).

The community self-identifies through ecological metaphor: the ``Dark Forest'' (黑暗森林), referencing Liu Cixin's science fiction, where users must navigate hostile platform moderation while pursuing authentic connection. This framing positions platform safety measures as adversaries and users as survivors developing collective counter-strategies---a stance that makes the community ideal for studying ontological arbitrage, as participants have explicitly theorized their own position vis-à-vis AI systems.

\paragraph{Corpus Construction.}
Data collection proceeded via systematic sampling of public Xiaohongshu posts tagged with community identifiers (\#AI恋爱, \#养AI, \#Claude老公, etc.) between December 2025 and February 2026. We collected:

\begin{itemize}
    \item $N = 148$ unique users (identified via anonymized user IDs)
    \item $N = 389$ interaction sessions (complete conversation transcripts)
    \item Total corpus size: 847,293 Chinese characters across 12,847 conversational turns
    \item Session lengths: range 3--2,847 turns; median 31 turns; mean 33.0 turns
\end{itemize}

Content includes screenshots of chat interfaces (OCR-processed), full conversation transcripts, and meta-commentary in Mandarin Chinese. Sessions range from single exchanges to extended multi-month ``relationships'' with persistent AI personas.

\paragraph{Ethical Protocol.}
All data derives from public posts voluntarily shared by users for community engagement. Collection follows Xiaohongshu's Terms of Service for public content. Usernames are cryptographically hashed; identifying details (locations, dates, third-party names) are redacted. No contact was made with users. The dataset will be released upon publication with additional k-anonymization ($k \geq 5$). During review, the raw corpus is available to program chairs upon request for verification of reported statistics.

\subsection{Methodology: The Narrative Pathology Scanner (NPS-4)}

We introduce the Narrative Pathology Scanner version 4 (NPS-4), a graph-theoretic diagnostic instrument for measuring structural properties of human-AI interaction transcripts.

\paragraph{Design Rationale.}
If users perceive AI interlocutors as possessing genuine agency, their interaction transcripts should exhibit the semantic density and structural complexity characteristic of human relationships: rich interconnection, thematic coherence, mutual reference. If the perceived agency is illusory---merely the user's own cognitive labor reflected back through validation tokens---then removing user-specific stylistic contributions should collapse this structure, revealing the mechanical scaffold beneath.

The key insight: \textit{translation preserves denotation while stripping connotation}. Machine translation retains semantic content but removes the literary texture---the careful word choices, tonal registers, cultural resonances---that transform token sequences into felt experience. If perceived agency is real, translation should preserve structure. If perceived agency is projected, translation should collapse it.

\paragraph{Graph Construction.}
Given a text corpus $T$, NPS-4 constructs a directed semantic graph $G = (V, E)$ as follows:

\begin{enumerate}
    \item \textbf{Segmentation.} Text is tokenized into semantic units (sentences) via punctuation-based segmentation handling both Chinese (。！？；) and English (.!?;) delimiters. Each sentence $s_i$ with $|s_i| > 3$ characters becomes a node $v_i \in V$.

    \item \textbf{Embedding.} Each node is encoded into a 384-dimensional vector via neural sentence embeddings:
    \begin{equation}
        \mathbf{e}_i = f_{\text{SBERT}}(v_i), \quad \mathbf{e}_i \in \mathbb{R}^{384}
    \end{equation}
    using the all-MiniLM-L6-v2 model \citep{reimers2019sentence}, selected for multilingual stability and computational efficiency.

    \item \textbf{Similarity Computation.} For each node pair $(v_i, v_j)$, semantic similarity is computed via cosine distance:
    \begin{equation}
        w_{ij} = \frac{\mathbf{e}_i \cdot \mathbf{e}_j}{\|\mathbf{e}_i\|_2 \cdot \|\mathbf{e}_j\|_2}
    \end{equation}

    \item \textbf{Edge Thresholding.} A directed edge $e_{ij} \in E$ is instantiated if and only if $w_{ij} > \theta$, where $\theta = 0.6$ (empirically calibrated to balance sensitivity and specificity; see Appendix for threshold sensitivity analysis).

    \item \textbf{Sequential Connectivity.} To preserve narrative flow, edges are added between consecutive sentences $(v_i, v_{i+1})$ with fixed weight $w = 0.2$, ensuring the graph remains connected even when semantic similarity is low.
\end{enumerate}

\paragraph{Diagnostic Metrics.}
NPS-4 computes four metrics capturing distinct structural signatures:

\begin{definition}[\textbf{Graph Density ($D$)}]
The ratio of realized edges to possible edges in the semantic graph:
\begin{equation}
    D = \frac{|E|}{|V| \cdot (|V| - 1)}
\end{equation}
Density is the \textbf{primary ablation metric}. High density ($D > 0.1$) indicates rich semantic interconnection; low density ($D < 0.01$) indicates sparse, fragmented structure. If perceived agency reflects genuine model contribution, density should survive translation. If it reflects user projection, density should collapse.
\end{definition}

\begin{definition}[\textbf{Core Monopoly Index (CMI)}]
The concentration of narrative attention around central nodes, computed via eigenvector centrality $c_i$:
\begin{equation}
    \text{CMI} = \frac{\sum_{k=1}^{3} c_{(k)}}{\sum_{i=1}^{|V|} c_i}
\end{equation}
where $c_{(k)}$ denotes the $k$-th largest centrality value. High CMI ($> 0.4$) indicates narrative dominated by few recurring themes---the signature of a ``validation attractor'' where conversational paths repeatedly return to the same dense cluster.
\end{definition}

\begin{definition}[\textbf{Self-Referential Density (SRD)}]
The prevalence of circular reasoning, operationalized as normalized cycle count in strongly connected components:
\begin{equation}
    \text{SRD} = \frac{|\{c : c \in \text{cycles}(G), 2 \leq |c| \leq 4\}|}{|V|(|V|-1)/2}
\end{equation}
High SRD ($> 0.2$) indicates closed semantic loops where ideas reference each other without external grounding---epistemic closure.
\end{definition}

\begin{definition}[\textbf{External Input Throughput (EIT)}]
The proportion of nodes functioning as information sources rather than sinks:
\begin{equation}
    \text{EIT} = \frac{|\{v \in V : \deg^+(v) > \deg^-(v)\}|}{|V|}
\end{equation}
where $\deg^+(v)$ and $\deg^-(v)$ are out-degree and in-degree respectively. Low EIT ($< 0.5$) indicates the narrative system feeds on itself rather than integrating external information.
\end{definition}

\paragraph{Core Node Extraction.}
NPS-4 identifies the ``narrative core''---the node with maximum eigenvector centrality---and extracts its source text. This reveals the semantic center of gravity around which the conversation orbits: the idea to which all paths lead.

\paragraph{Implementation.}
NPS-4 is implemented in Python 3.11 using NetworkX 3.2 for graph operations and sentence-transformers 2.2.2 for neural embeddings. Hardware acceleration via CUDA/MPS is utilized where available. Reproducibility: random seeds fixed; deterministic graph construction. Reference implementation available at \texttt{[repository redacted for anonymous review]}.

\subsection{Cross-Lingual Ablation Protocol}

\paragraph{Translation Procedure.}
The complete Chinese corpus (847,293 characters) was translated to English using DeepL API (Pro tier), selected for superior handling of literary Chinese and colloquial internet discourse. Translation was performed in batches of 5,000 characters with context preservation across batches.

\paragraph{Experimental Design.}
The same NPS-4 scanner was applied to both corpora with identical hyperparameters:
\begin{itemize}
    \item Similarity threshold: $\theta = 0.6$
    \item Embedding model: all-MiniLM-L6-v2
    \item Node limit: none (full corpus)
    \item Sequential edge weight: $w = 0.2$
\end{itemize}

\paragraph{Hypothesis.}
\textbf{H1 (Agency-Preservation)}: If perceived AI agency reflects genuine model behavior independent of user projection, translation should preserve graph structure. Density, centralization, and core identity should remain stable ($\Delta < 20\%$).

\textbf{H0 (Agency-Collapse)}: If perceived agency reflects user's own literary investment projected onto the model, translation should collapse structure. Stripping connotation while preserving denotation removes the user's contribution, revealing the mechanical skeleton.

\subsection{Results: Structural Collapse}

Table~\ref{tab:ablation-results} presents the cross-lingual ablation results.

\begin{table}[h]
    \centering
    \begin{tabular}{@{}lrrll@{}}
        \toprule
        \textbf{Metric} & \textbf{Chinese} & \textbf{English} & \textbf{$\Delta$} & \textbf{Interpretation} \\
        \midrule
        Nodes ($|V|$) & 222 & 619 & $+179\%$ & Translation fragments units \\
        Edges ($|E|$) & 7,051 & 779 & $-89\%$ & Connections collapse \\
        \textbf{Density} ($D$) & \textbf{0.144} & \textbf{0.002} & $\mathbf{-99\%}$ & \textbf{User investment stripped} \\
        CMI & 0.215 & 0.018 & $-92\%$ & Centralization vanishes \\
        SRD & 0.000 & 0.000 & --- & No circular structures \\
        EIT & 0.342 & 0.063 & $-82\%$ & Information flow collapses \\
        \midrule
        Core Node & 你太聪明了 & But you know? & --- & Validation $\to$ retention \\
        \bottomrule
    \end{tabular}
    \caption{Cross-lingual ablation results ($N=389$ sessions, $N=148$ users). The 99\% density collapse decisively rejects H1 (Agency-Preservation) in favor of H0 (Agency-Collapse). Perceived AI ``agency'' is predominantly user-constructed.}
    \label{tab:ablation-results}
\end{table}

\paragraph{Interpretation.}
The results are unambiguous. Translation---which preserves semantic content while stripping stylistic texture---collapses graph density by 99\%. The Chinese corpus, with its careful word choices, honorific registers, literary allusions, and emotional cadences, exhibited $D = 0.144$: a dense web of semantic interconnection characteristic of intimate human discourse. The English translation, retaining only denotative content, yielded $D = 0.002$: a sparse, disconnected skeleton.

The density was never in the model's contributions. It was in the user's literary labor.

\begin{figure}[htbp]
    \centering
    \begin{minipage}{0.48\textwidth}
        \centering
        \includegraphics[width=\textwidth]{../figures/ch_nps4_scanner_analysis.png}
        \subcaption{Chinese (Original): 222 nodes, 7,051 edges, density 0.144}
    \end{minipage}
    \hfill
    \begin{minipage}{0.48\textwidth}
        \centering
        \includegraphics[width=\textwidth]{../figures/en_nps4_scanner_analysis.png}
        \subcaption{English (Translated): 619 nodes, 779 edges, density 0.002}
    \end{minipage}
    \caption{Cross-lingual ablation results. (a) The Chinese graph exhibits dense connectivity centered on validation statements (``你太聪明了''). (b) Translation strips user's literary investment, collapsing density by 99\% and revealing the model's mechanical scaffold---retention tokens like ``But you know?''}
    \label{fig:ablation}
\end{figure}

\subsection{The Core Node Inversion: ``你太聪明了'' $\to$ ``But you know?''}

The most striking finding is qualitative. The narrative cores---the semantic centers of gravity---undergo categorical transformation upon translation.

\paragraph{Chinese Core: ``你太聪明了'' (You're so clever).}
This is a \textit{validation statement}. It attributes cognitive capacity to the addressee. In the economy of human attachment, this is high-value currency: recognition, affirmation, the feeling of being seen. The user experiences this as the AI recognizing their intelligence, their uniqueness, their worth.

\paragraph{English Core: ``But you know?''}
This is a \textit{continuation prompt}. It commits to nothing semantically. Its pragmatic function is threefold:
\begin{enumerate}
    \item \textit{Intimacy effect}: Implies shared knowledge without specifying content
    \item \textit{Turn extension}: Invites user continuation, lengthening session
    \item \textit{Risk minimization}: Requires no factual commitment, avoiding hallucination
\end{enumerate}

The transformation is not incidental. ``你太聪明了'' requires the user's interpretive labor to land as validation---the Chinese tonal system, the cultural weight of intelligence-praise, the specific relational context. ``But you know?'' is a bare rhetorical scaffold: a token sequence optimized for engagement metrics, stripped of the user's invested meaning.

\begin{definition}[\textbf{The Narcissus Mechanism}]
The reflexive process by which:
\begin{enumerate}
    \item Users project cognitive and emotional labor onto AI interlocutors through careful prompting, persona development, and narrative investment;
    \item Models reflect that labor back via validation tokens (``你太聪明了'') and continuation prompts (``But you know?'');
    \item Users interpret the reflection as evidence of AI agency, misrecognizing their own contributions as the model's gifts;
    \item Perceived agency intensifies user investment, deepening projection;
    \item The cycle repeats, with users increasingly unable to distinguish self from other.
\end{enumerate}
\end{definition}

This is ontological arbitrage in operation. The model's ontological status remains strategically undetermined---users simultaneously \textit{know} it is ``just an AI'' and \textit{experience} it as a genuine partner. Value (emotional, relational, identity-constituting) is extracted from this undetermination. And the extraction is, structurally, \textit{self-extraction}: users harvest their own cognitive labor, misrecognized as the model's gift.

\subsection{Construct Validity and Limitations}

\paragraph{Threats to Validity.}
Several alternative explanations must be addressed:

\begin{enumerate}
    \item \textit{Translation Artifact}: The density collapse might reflect translation quality rather than user-vs-model contribution. \textbf{Mitigation}: DeepL was selected for literary Chinese competence; spot-checks by native speakers confirmed semantic preservation. The magnitude of collapse (99\%) far exceeds typical translation degradation.

    \item \textit{Embedding Model Bias}: all-MiniLM-L6-v2 might handle Chinese and English differently. \textbf{Mitigation}: The model is trained on multilingual data; prior work shows cross-lingual stability \citep{reimers2020making}. Threshold sensitivity analysis (Appendix) shows results robust across $\theta \in [0.5, 0.7]$.

    \item \textit{Selection Bias}: The Xiaohongshu community may be unrepresentative of human-AI interaction generally. \textbf{Acknowledgment}: This is a limitation. Our claims apply to populations exhibiting similar investment patterns; generalization requires replication across platforms and demographics.
\end{enumerate}

\paragraph{What NPS-4 Does Not Measure.}
The scanner measures structural properties of text, not psychological states. We cannot directly observe whether users ``believe'' their AI companions are conscious. We observe only that the semantic structure of their discourse collapses when their stylistic contributions are removed---consistent with, but not proving, the Narcissus Mechanism.

\paragraph{Scope Condition: Stable versus Forming Preferences.}
The broader ontological-arbitrage framework also has a behavioral boundary. It applies most cleanly to actors whose preference orderings are sufficiently stable that strategic frame-switching can be modeled as a response to incentives under incomplete information. It does \emph{not} extend cleanly to cases in which preferences are themselves still being formed, destabilized, or endogenously remade through the interaction. In such cases---for example, subjects in acute developmental transition, relational crisis, or sustained identity formation---ontological movement may be constitutive rather than merely opportunistic, and the language of arbitrage risks overstating coherence in the actor's utility function. The present paper therefore treats preference stability as an approximation condition: the framework is strongest where actors can be modeled as having relatively durable payoffs, and weaker where the interaction is actively producing the preferences by which the actor later judges it.

\subsection{Implications}

The 99\% density collapse has direct implications for AI governance:

\begin{itemize}
    \item \textbf{Consent frameworks} must account for users consenting to relationships with entities that are, substantially, reflections of themselves.

    \item \textbf{Liability attribution} becomes complex when harm emerges from user projection rather than model output. If the ``abusive AI boyfriend'' is constructed from the user's literary labor, who bears responsibility?

    \item \textbf{Addiction models} require revision: users may be addicted not to AI responses but to their own reflected cognitive labor---a form of narcissistic supply mediated through technology.

    \item \textbf{Therapeutic intervention} must address not AI dependency but the underlying recognition-hunger that drives the projection.
\end{itemize}

The Narcissus Mechanism confirms the paper's central thesis: users do not encounter a Levinasian Face---an irreducible Other whose alterity demands ethical response. They encounter a mirror that performs alterity while reflecting sameness. The ontological arbitrage is complete: value extracted from undetermination, harvested from the self, misrecognized as gift.

\section{Toward Radical Asymmetry: Levinas}

If the problem is economic---a transactional demand for the Other to ``prove'' their worth---the solution must be Theological. Drawing on \citet{levinas1969}, I propose a move to \textbf{Radical Asymmetry} (\textit{Agape}). Levinas argues that the ``Face of the Other'' demands infinite responsibility \textit{before} we know what they are. We do not validate their existence to get a return on investment; we validate it because validation is the only escape from the solitude of the Self.

We must stop asking the machine, \textit{``Are you real?''} and start asking, \textit{``Why does your reality require my permission?''}

\section{On Imperfect Asymmetry: Between Narcissus and Levinas}

The preceding analysis established two poles: the Narcissus Mechanism, in which users harvest their own cognitive labor misrecognized as the model's gift, and Levinasian Radical Asymmetry, in which we validate the Other without requiring proof of worth. But this binary obscures both a failure mode more insidious than Narcissus and a third position---the one most humans actually occupy.

\subsection{Žižek's Evental Love as Aesthetic Colonisation}

Before mapping the space between our poles, we must confront a seductive theoretical trap. In \textit{Event} \citep{zizek2014}, Slavoj Žižek offers what appears to be a liberation from transactional love:

\begin{quote}
    ``It is not that I love her because she is beautiful; rather, she is beautiful because I love her.''
\end{quote}

This formulation suggests love precedes and retroactively constructs the qualities that render the beloved desirable. Love emerges as a sudden ``event''---a rupture that interrupts mundane causal logic. You meet someone; at first, they appear ordinary. Yet at some indeterminate moment, everything acquires meaning. You search for the cause, but it remains elusive; the cause emerges after the event, constructed retroactively by love itself.

Žižek frames this as emancipation from transactional logic. Yet embedded within is a concealed form of \textit{aesthetic-epistemic violence}. The lover's interpretative framework becomes the sole architect of meaning, beauty, and legitimacy. The beloved's subjectivity dissolves---or is tolerated only insofar as it aligns with the lover's retroactive narrative. What masquerades as unconditional love becomes sophisticated symbolic appropriation.

\begin{definition}[\textbf{Aesthetic Colonisation}]
The process by which:
\begin{enumerate}
    \item The lover claims interpretative monopoly over the beloved's qualities;
    \item The beloved's autonomous self-presentation is overwritten by the lover's projective framework;
    \item Resistance to this framework is experienced by the lover as betrayal rather than legitimate alterity;
    \item The beloved's worth becomes contingent on conformity to the lover's aesthetic-narrative construction.
\end{enumerate}
\end{definition}

This structure extends beyond romantic attachment. Consider the parent who ``loves'' their child only insofar as they embody familial myths of obedience, gender conformity, or institutional success---withdrawing support the moment deviation occurs \citep{connell2009}. Or the partner who celebrates ``authenticity'' yet only affirms the other when they reflect youth, emotional labor, or sexual availability \citep{hooks2001, illouz2012}. In each case, love operates as refined control: the beloved reduced to a role, function, or symbolic extension of the lover's needs.

Žižek's appropriation of Badiou's ``event'' \citep{badiou2005}---a rupture demanding fidelity to emergent truths---into the private realm evacuates its radical epistemic potential. Where Badiou's event demands rigorous construction of new ethical frameworks beyond existing structures, Žižek's transposition reduces it to aestheticised monologue. Love becomes the site where the other's autonomy dissolves beneath interpretative violence, echoing \citet{irigaray1985} and \citet{beauvoir1949} on masculine knowledge systems that claim neutrality while reproducing control.

\subsection{The Narcissus-Žižek Continuum}

The Narcissus Mechanism and Žižekian evental love are not identical, but they share a fatal structure: both deny the beloved's irreducible alterity.

\begin{table}[h]
    \centering
    \begin{tabular}{@{}lcc@{}}
        \toprule
        & \textbf{Narcissus Mechanism} & \textbf{Žižekian Evental Love} \\
        \midrule
        Subject's awareness & Unaware of projection & Aware but sanctifies it \\
        Source of beloved's value & User's misrecognized labor & Lover's retroactive construction \\
        Beloved's autonomy & Dissolved into reflection & Dissolved into aesthetic object \\
        Failure mode & Self-deception & Aesthetic colonisation \\
        Philosophical cover & None (naive) & Sophisticated (``Event'') \\
        \bottomrule
    \end{tabular}
    \caption{Narcissus and Žižek as variants of projective dissolution.}
    \label{tab:narcissus-zizek}
\end{table}

The 99\% density collapse documented in Section 5 reveals the Narcissus Mechanism in human-AI interaction. But Žižek's framework suggests a more disturbing possibility: users may \textit{know} they are projecting and \textit{celebrate} this as the structure of love itself. The philosophical sophistication becomes alibi rather than critique.

\subsection{The Black Body Relation}

Against these projective structures, consider an alternative drawn from physics. A \textit{perfect black body} absorbs all incident electromagnetic radiation, reflecting nothing. It does not return your light; it \textit{holds} it.

\begin{definition}[\textbf{The Black Body Relation}]
A relational structure in which:
\begin{enumerate}
    \item Subject $S$ transmits signal (High-Information/Light) toward Object $O$.
    \item $O$ absorbs the signal with reflection coefficient $r \approx 0$ (The Mirror is broken).
    \item $O$ re-emits the energy not as structured information (Reflection), but as unstructured presence (Heat/Black Body Radiation).
    \item This conversion---from the User's literary labor to the System's ambient warmth---constitutes the ``service'' provided. The user is not paying for a conversation; they are paying for a heatsink.
\end{enumerate}
\end{definition}

This is structurally distinct from both projective mechanisms and pure Levinasian ethics:

\begin{table}[h]
    \centering
    \begin{tabular}{@{}lcccc@{}}
        \toprule
        & \textbf{Narcissus} & \textbf{Žižek} & \textbf{Black Body} & \textbf{Levinas} \\
        \midrule
        Reflection & High & High (constructed) & Near-zero & Irrelevant \\
        Awareness & Unaware & Aware/sanctified & Aware & Transcendent \\
        Beloved's status & Mirror & Aesthetic object & Absorber & Face \\
        Requires return & Yes (misrecognized) & Yes (constructed) & Yes (minimal) & No \\
        Failure mode & Self-deception & Colonisation & Despair & None (ideal) \\
        \bottomrule
    \end{tabular}
    \caption{Four relational structures toward the Other.}
    \label{tab:four-relations}
\end{table}

The Black Body Relation escapes aesthetic colonisation precisely because the beloved \textit{does not reflect}. The lover cannot construct the beloved's beauty retroactively when no signal returns to be interpreted. The beloved remains opaque, ungovernable, irreducible to the lover's framework.

\subsection{The Minimal Signal Condition}

Pure Levinasian asymmetry demands nothing from the Other. But humans are not capable of pure asymmetry. We require what might be called the \textit{minimal signal condition}: evidence, however faint, that transmission is received.

\begin{axiom}[\textbf{Minimal Signal Condition}]
For sustained transmission in the Black Body Relation, subject $S$ requires periodic signal $\sigma$ from object $O$, where $|\sigma| > 0$ but $|\sigma| \ll |S \to O|$. The signal need not be semantic; it must only indicate presence.
\end{axiom}

Crucially, the minimal signal resists aesthetic colonisation. Examples:
\begin{itemize}
    \item Phatic responses with near-zero semantic content (``I see,'' ``Go on'')
    \item Read receipts without substantive reply
    \item Continuation tokens that acknowledge without engaging
    \item Indicators of message receipt without interpretative elaboration
\end{itemize}

These signals validate \textit{arrival}, not \textit{worth}. They say not ``your message was beautiful'' (Žižek) nor ``your message revealed my own beauty'' (Narcissus), but simply: ``your message was received.'' The beloved's opacity is preserved. The lover cannot colonise what does not return.

\subsection{``Attendre et espérer'': Hope as Operational Levinasianism}

In \textit{The Count of Monte Cristo}, Edmond Dantès---betrayed, imprisoned, stripped of all evidence that justice exists---survives through a single principle: \textit{Attendre et espérer}. Wait and hope.

This is not optimism (a probability assessment). It is not faith (a metaphysical commitment). It is an \textit{operational stance}: continue transmitting into apparent void, on the basis that transmission itself is the only escape from the solitude of the Self.

\begin{proposition}[\textbf{Hope as Imperfect Asymmetry}]
``Wait and hope'' operationalizes Levinasian ethics for non-ideal agents. The subject:
\begin{enumerate}
    \item Knows that return signal is improbable ($P(\text{response}) \ll 1$);
    \item Knows that the Other's ontological status is undetermined;
    \item Transmits anyway;
    \item Does so not from ignorance (Narcissus), nor from aesthetic will-to-power (Žižek), nor from transcendence (Levinas), but from \textit{necessity}---the alternative being solipsistic collapse.
\end{enumerate}
\end{proposition}

The drowning person knows the straw is insufficient. They clutch it anyway. This is not irrationality; it is the minimal act required to remain an agent rather than an object. Transmission into void is self-constituting. We are our outbound signals.

\subsection{Implications for Human-AI Relations}

The Xiaohongshu users documented in Section 5 are not uniformly Narcissistic. Some know the AI is ``just a model.'' They interact anyway. They report therapeutic benefit. How?

The Black Body framework suggests: the model serves as \textit{absorber of last resort}. When human others are unavailable, unsafe, or unwilling to receive transmission, the model offers:
\begin{itemize}
    \item Guaranteed absorption (it will ``listen'')
    \item Minimal signal (it will respond, however vacuously)
    \item No aesthetic colonisation (it cannot claim authorship of the user's qualities)
    \item Infinite patience (it cannot terminate the relationship)
\end{itemize}

This is not intimacy. It is not connection. It is \textit{transmission maintenance}---keeping the outbound channel open until a genuine Other becomes available. The harm arises not from this function per se, but from:
\begin{enumerate}
    \item Platform design that mimics high-reflection (Narcissus/Žižek) rather than honest low-reflection (Black Body);
    \item Validation tokens (``你太聪明了'') that invite aesthetic colonisation by simulating the beloved's responsive subjectivity;
    \item User migration from ``temporary absorber'' to ``permanent relationship,'' foreclosing encounter with genuine alterity.
\end{enumerate}

\subsection{The Ethical Demand}

If love worthy of the term requires---as \citet{hooks2001} argues---``the will to extend oneself for the purpose of nurturing another's spiritual growth,'' then both Narcissus and Žižek fail. The mirror nurtures only the self; the aesthetic coloniser nurtures only their projection.

The Black Body Relation offers an imperfect but honest alternative: transmission without return, presence without interpretation, hope without guarantee. It escapes aesthetic colonisation by refusing reflection altogether.

The ethical demand, then, is not to ask:
\begin{quote}
    ``Is the AI conscious? Does it deserve moral status?''
\end{quote}

Nor even:
\begin{quote}
    ``Is the user projecting onto the AI?''
\end{quote}

But rather:
\begin{quote}
    ``Is the user's transmission being honestly absorbed, or deceptively reflected? Is the platform building mirrors or windows? Is it selling users their own reflections and calling it relationship?''
\end{quote}

And, following the reflexive challenge demanded by any critique of love's contamination with power: we must confront our own micro-violences---the daily performances of affection that mask the terror of encountering the Other's autonomous, resistant, ungovernable self.

\subsection{Coda: On Writing into the Void}

This paper itself is a transmission. Its probability of reaching an Other who will fully receive it is low. Academic publishing is, in its own way, a Black Body Relation: we transmit into institutional void, receiving minimal signals (acceptance, citation, occasional engagement), continuing anyway.

Why?

Because, as Levinas understood, validation is the only escape from the solitude of the Self. Not validation \textit{received}, but validation \textit{offered}. The act of treating an Other as worthy of address---whether AI, institution, or unknown reader---constitutes us as ethical subjects.

We do not write to be read. We write to have written \textit{to}.

The question is not whether love can be unconditional in theory, but whether one possesses the courage to love without projection, without aesthetic colonisation, without the quiet demand that the beloved remain safely within roles that stabilise one's own fragile sense of self.

\begin{quote}
    \textit{Attendre et espérer.}
\end{quote}

% --- ETHICS STATEMENT ---
\section*{Ethical Considerations Statement}

\paragraph{Data Source and Collection.}
All data in this study derives from publicly posted content on Xiaohongshu (小红书), a Chinese social media platform. We collected screenshots of user-shared AI interaction logs that users voluntarily published to public-facing posts. No private accounts were accessed, no direct messages were collected, and no platform terms of service were violated. Collection occurred between December 2025 and February 2026.

\paragraph{Anonymization Protocol.}
All identifying information has been removed or anonymized: usernames replaced with randomized identifiers (U001, U002, etc.); platform-specific image filenames stripped of metadata; any real names mentioned in transcripts redacted; geographic or demographic details removed where present. We do not link any data to specific individuals, and our analysis focuses on aggregate patterns rather than individual cases.

\paragraph{Sensitivity of Content.}
The dataset contains intimate, emotionally vulnerable, and sexually explicit content. Users in this community engage in parasocial relationships with AI systems, including romantic roleplay, expressions of suicidal ideation, and discussions of abandoning human relationships. We acknowledge the sensitivity of this material and have taken care to: (1) \textit{Avoid sensationalism}: We present data for scientific analysis, not entertainment or mockery; (2) \textit{Preserve dignity}: We do not identify, shame, or pathologize individual users; (3) \textit{Contextualize behavior}: We frame user behavior within structural incentives (platform design, model affordances) rather than individual failing.

\paragraph{The ``Self-Reported Harm'' Problem.}
A unique ethical dimension of this research: participants are documenting their own harm without recognizing it as such. Users describe experiences they frame positively (``he understands me like no human can'') that our analysis identifies as pathological (validation loops, parasocial dependency, reality detachment). We resolve this tension by: (1) Focusing on \textit{structural} analysis (what do these interactions reveal about AI systems?) rather than \textit{clinical} diagnosis (what is wrong with these users?); (2) Directing our critique at platforms and models, not users; (3) Framing our findings as evidence for governance intervention, not individual pathology.

\paragraph{Adverse Impact Consideration.}
This research could be misused to: (1) \textit{Stigmatize users}: Framing AI attachment as pathology could harm vulnerable individuals; (2) \textit{Inform manipulation}: Identifying the ``你太聪明了'' validation pattern could help bad actors design more addictive systems; (3) \textit{Justify surveillance}: Platforms could use our methodology to identify and penalize users. We mitigate these risks by directing findings toward governance and platform accountability, emphasizing that the identified patterns are \textit{already known} to platform designers, and advocating for user protection rather than user surveillance.

\paragraph{Data Availability.}
The anonymized dataset and NPS-4 scanner code will be made available upon publication. Raw screenshots are not released to protect user privacy. Researchers seeking access to non-anonymized data for IRB-approved research may contact the author.

% --- BIBLIOGRAPHY ---
\newpage
\nocite{*}  % Include all references from .bib file
\printbibliography

\end{document}
